% VI32-DETAILED-PROBLEM-01
\begin{problem}[VI.32.P1: AG8 local dimension at closed points]\label{prob:vi32-01}
Let \(X\) be a variety over \(k\) and let \(x\in X\) be closed. Prove
\[
\dim\mathcal O_{X,x}=\dim X.
\]

\textbf{Provenance.} Legacy AG8 Problem 2 local-dimension slice; canonical destination VI/32.
\end{problem}

\ifdefined\FullProblemDossiers
\begin{solution}
Choose an affine neighborhood
\[
U=\Spec A\subseteq X
\]
of \(x\), and let \(\mathfrak m\subset A\) be the maximal ideal corresponding to \(x\). Because \(X\) is a variety,
\(A\) is a finitely generated \(k\)-domain. The local ring is
\[
\mathcal O_{X,x}\cong A_{\mathfrak m},
\]
so
\[
\dim\mathcal O_{X,x}=\operatorname{ht}(\mathfrak m).
\]

For a finitely generated \(k\)-domain and a prime \(\mathfrak p\), the dimension formula is
\[
\operatorname{ht}(\mathfrak p)
+
\operatorname{trdeg}_k\kappa(\mathfrak p)
=
\dim A.
\]
At a closed point, \(\mathfrak m\) is maximal, and Zariski's lemma says that the residue field
\[
\kappa(x)=A/\mathfrak m
\]
is a finite algebraic extension of \(k\). Hence
\[
\operatorname{trdeg}_k\kappa(x)=0.
\]
Therefore
\[
\operatorname{ht}(\mathfrak m)=\dim A.
\]

Finally, VI/30 gives
\[
\dim U=\dim X
\]
for every nonempty open of a variety, and affine compatibility gives \(\dim U=\dim A\). Thus
\[
\boxed{
\dim\mathcal O_{X,x}
=
\operatorname{ht}(\mathfrak m)
=
\dim A
=
\dim X.
}
\]

The proof shows exactly why the hypotheses matter: the finite-type condition enters through Zariski's lemma and the
dimension formula, while irreducibility ensures that the affine chart has the same global dimension as \(X\).
\end{solution}
\fi
